\section{Contenido del Trabajo}

\lipsum[4]

Enla cita dentro del texto~\textcite{Wolf2017} y al final del texto~\parencite{Margolis2020}...

\subsection{Punto 1}

\lipsum[5]

En la Figura~\ref{fig:punto1} se observa ....

\insertarFigura{punto1.png}{1.0}{Imagen de ejemplo Punto 1.}{fig:punto1}

\lipsum[6]

En la Tabla~\ref{tab:voltajes} se muestra....

\begin{table}[]
    \centering
    \caption{Resultados de la medición de voltajes en el microcontrolador.}
    \label{tab:voltajes}
    \begin{tabular}{@{} lcc @{}} % l=izq, c=centro, c=centro
        \toprule
        \textbf{Componente} & \textbf{Volt. Nominal (V)} & \textbf{Volt. Medido (V)} \\ \midrule
        Puerto A            & 5.0                          & 4.92                        \\
        Puerto B            & 3.3                          & 3.28                        \\
        VCC                 & 5.0                          & 5.01                        \\ \bottomrule
    \end{tabular}
\end{table}

\lipsum[7]

\begin{table}[ht] 
\caption{Performances as amplifier and rectifier at 10.1 GHz} %title of the table 
\centering % centering table 
\begin{tabular}{c rr  rr  } % creating eight columns 
%\hline\hline %inserting double-line 
\toprule
 &\multicolumn{2}{c}{Amplifier} & \multicolumn{2}{c}{Rectifier}\\  
Circuit &A&B&A&B\\ [0.25ex] 
%\hline
\midrule
Max Efficiency(\%) & 67.87 & 56.47 & 64.40 & 63.94 \\  
DC Power (mW) & 4186.0 & 5112.0 & 1671.0 & 3182.0\\ 
RF Power (mW) & 3281.0$^{\mathrm{a}}$ & 3362.0 & 2594.0 & 4976.0\\  
%\hline\hline
\bottomrule
\multicolumn{5}{l}{$^{\mathrm{a}}$Sample of a Table footnote.}
\end{tabular} 
\label{tab:hresult} 
\end{table} 


En el Código~\ref{cod:python_ejemplo} se muesta...

\lstinputlisting[
    style=python_style,
    linerange={11-},
    firstnumber=11,
    caption={Implementación del algoritmo de control en Python.},
    label={cod:python_ejemplo}
]{src/punto1.py}

\lipsum